% PatchMatch一维教学例：A中位置0-3的正确偏移均为+10（匹配B中10-13）。
% 初始化(+3,+10,-1,+5)；第一次正向传播后(+3,+10,+10,+10)；第二次逆向传播后全部为+10。
% 三个面板上下排列，每格5mm，宽约107mm、高约108mm。
\begin{tikzpicture}[x=5mm,y=1mm,font=\figurefont\small,
 cellA/.style={draw=BookBlue,fill=BookBlue!10,line width=.6pt},
 cellB/.style={draw=BookMuted,fill=white,line width=.5pt},
 good/.style={-{Stealth[length=1.6mm]},BookTeal,line width=.9pt},
 bad/.style={-{Stealth[length=1.6mm]},BookBlue,dashed,line width=.8pt},
 prop/.style={-{Stealth[length=1.4mm]},BookAmber,line width=.8pt}]
 \foreach \panel/\ytop/\title/\offs in {
   0/0/{初始化}/{3,10,-1,5},
   1/-34/{第一次迭代：正向传播}/{3,10,10,10},
   2/-68/{第二次迭代：逆向传播}/{10,10,10,10}} {
  \node[anchor=west] at (-4.6,\ytop+5) {\title};
  \node[anchor=east] at (-.2,\ytop-3) {$A$};
  \node[anchor=east] at (-.2,\ytop-19) {$B$};
  \foreach \i in {0,...,3} {
   \draw[cellA] (\i,\ytop-6) rectangle (\i+1,\ytop);
   \node at (\i+.5,\ytop-3) {\footnotesize\i};
  }
  \foreach \j in {0,...,15} {
   \draw[cellB] (\j,\ytop-22) rectangle (\j+1,\ytop-16);
   \node[text=BookMuted] at (\j+.5,\ytop-19) {\scriptsize\j};
  }
  \foreach \o [count=\k from 0] in \offs {
   \pgfmathtruncatemacro{\tgt}{\k+\o}
   \ifnum\o=10
    \draw[good] (\k+.5,\ytop-6.3) -- (\tgt+.5,\ytop-15.7);
   \else
    \draw[bad] (\k+.5,\ytop-6.3) -- (\tgt+.5,\ytop-15.7);
   \fi
  }
 }
 \draw[prop] (1.6,-34+0.8) to[bend left=50] (2.4,-34+0.8);
 \draw[prop] (2.6,-34+0.8) to[bend left=50] (3.4,-34+0.8);
 \draw[prop] (1.4,-68+0.8) to[bend right=50] (0.6,-68+0.8);
 \draw[good] (3,-96)--(4,-96); \node[anchor=west] at (4.1,-96) {正确偏移};
 \draw[bad] (7.5,-96)--(8.5,-96); \node[anchor=west] at (8.6,-96) {错误偏移};
 \draw[prop] (12,-96) to[bend left=40] (13,-96); \node[anchor=west] at (13.1,-96) {传播方向};
\end{tikzpicture}
